\begin{frame}{제어 비용 가중치 1.0: 보상 함수는 공학 제재의 번역}
  제어 비용의 가중치가 \textbf{1.0}이라는 점에 주의 --- docstring에는
  0.1이라고 적혀 있으나, 소스의 기본값은
  \texttt{reward\_control\_weight = 1}.
  \vskip0.4em
  \begin{center}
    \kq{문서와 코드가 다르면 \textbf{코드가 맞다} --- 이 분야에서 자주
    일어나는 일.}
  \end{center}
  \vskip0.4em
  ``행동이 크면 벌을 주는'' 항을 왜 넣나? 실물 로봇에서 큰 토크는
  \textbf{모터의 부담 · 발열 · 기계 마모}이기 때문.
  \begin{block}{보상 함수는 공학 제재가 수학으로 번역되는 장소}
    이 항을 무시한 정책은 시뮬레이션에서 ``자주 흔들어서 근접''하는
    전략을 배울 수 있지만, 실물에서는 \textbf{액추에이터를 고치는}
    정책.
  \end{block}
\end{frame}
